% ======================================================================
%  CHƯƠNG 1 — MỞ ĐẦU
%  Yêu cầu (mục 2.2, 2.6 của đề bài): bối cảnh & động lực; phát biểu bài toán
%  (Input/Output/Ràng buộc); mục tiêu & phạm vi; đóng góp; cấu trúc báo cáo.
%  Độ dài tham khảo: 5--10 trang.
% ======================================================================
\chapter{Mở đầu}
\label{ch:gioithieu}

\section{Bối cảnh và động lực nghiên cứu}
\label{sec:boicanh}
Sự phổ biến của học máy trong nhiều lĩnh vực đã kéo theo nhu cầu ngày càng lớn
về việc xây dựng mô hình một cách nhanh chóng và hiệu quả. Tuy nhiên, một quy
trình học máy điển hình đòi hỏi nhiều bước thủ công đòi hỏi chuyên môn cao: tiền
xử lý dữ liệu, trích chọn đặc trưng, lựa chọn thuật toán và tinh chỉnh siêu tham
số. Học máy tự động (\textit{Automated Machine Learning}, AutoML) ra đời nhằm tự
động hoá các bước này, giúp người dùng không chuyên vẫn có thể thu được mô hình
chất lượng với công sức tối thiểu~\cite{feurer2015autosklearn}.

Hiện nay có nhiều thư viện AutoML mã nguồn mở phổ biến như FLAML~\cite{wang2021flaml},
H2O~AutoML~\cite{ledell2020h2o} và AutoGluon~\cite{erickson2020autogluon}, mỗi thư
viện dựa trên một triết lý thiết kế khác nhau. Điều này đặt ra một câu hỏi thực
tiễn quan trọng: \textit{với một bài toán và một ngân sách tài nguyên cho trước,
nên lựa chọn thư viện AutoML nào?} Trả lời câu hỏi này đòi hỏi một quá trình đánh
giá so sánh khách quan.

Trên thực tế, việc so sánh các thư viện AutoML thường không công bằng và khó tái
lập. Các báo cáo khác nhau thường sử dụng ngân sách thời gian, cách chia dữ liệu,
độ đo và cấu hình tài nguyên khác nhau; tệ hơn, các lần chạy thất bại đôi khi bị
loại bỏ một cách thầm lặng, khiến kết quả trở nên thiên lệch. Nhằm khắc phục,
Gijsbers và cộng sự đề xuất \textit{AutoML Benchmark} (AMLB) --- một khung đánh
giá áp dụng \emph{cùng một giao thức thống nhất} (cùng cách chia fold, cùng ngân
sách thời gian, cùng độ đo, cùng tài nguyên) cho mọi thư viện và ghi nhận đầy đủ
các lần thất bại~\cite{gijsbers2024amlb}. Đây chính là động lực của đồ án: xây
dựng một quy trình đánh giá AutoML \emph{công bằng, tái lập được} và thuận tiện
sử dụng trên dữ liệu thực do người dùng lựa chọn.

\section{Phát biểu bài toán}
\label{sec:phatbieu}
Bài toán của đồ án được phát biểu như sau. Cho một tập các bộ dữ liệu
$\mathcal{D} = \{D_1, \dots, D_m\}$ (bao gồm cả bài toán phân lớp và hồi quy) và
một tập các thư viện AutoML $\mathcal{F} = \{F_1, \dots, F_k\}$. Với một ngân
sách tài nguyên thống nhất (thời gian, số nhân CPU, bộ nhớ), cần đánh giá và
\emph{xếp hạng} các thư viện trong $\mathcal{F}$ theo chất lượng dự đoán và theo
sự đánh đổi giữa chất lượng và chi phí, sao cho phép so sánh là công bằng và có
thể tái lập. Bảng~\ref{tab:io} tóm tắt bài toán theo ba thành phần Input --
Output -- Ràng buộc.

\begin{table}[H]
  \centering
  \small
  \caption{Phát biểu bài toán theo Input -- Output -- Ràng buộc.}
  \label{tab:io}
  \renewcommand{\arraystretch}{1.3}
  \begin{tabularx}{\linewidth}{@{}l X@{}}
    \toprule
    \textbf{Thành phần} & \textbf{Mô tả} \\
    \midrule
    \textbf{Input} &
    Tập bộ dữ liệu $\mathcal{D}$ dạng bảng (mỗi bộ gồm ma trận đặc trưng và cột
    nhãn/giá trị mục tiêu, có xác định loại tác vụ: phân lớp nhị phân, phân lớp
    đa lớp hoặc hồi quy); tập thư viện AutoML $\mathcal{F}$; ngân sách tài nguyên
    thống nhất (thời gian trên mỗi fold, số nhân CPU, bộ nhớ). \\
    \textbf{Output} &
    Với mỗi cặp (thư viện, bộ dữ liệu): điểm số theo độ đo phù hợp
    (AUC cho phân lớp nhị phân, log-loss cho phân lớp đa lớp, RMSE cho hồi quy)
    cùng chi phí thời gian và bộ nhớ; bảng xếp hạng tổng hợp các thư viện theo
    thứ hạng trung bình; biểu đồ đánh đổi giữa chất lượng và thời gian. \\
    \textbf{Ràng buộc} &
    Giao thức đánh giá phải \emph{đồng nhất} giữa các thư viện (cùng fold, cùng
    ngân sách, cùng độ đo, cùng tài nguyên); các lần chạy thất bại phải được
    \emph{ghi nhận} thay vì bỏ qua; toàn bộ quy trình phải tái lập được và chạy
    trong môi trường cô lập bằng Docker. \\
    \bottomrule
  \end{tabularx}
\end{table}

\section{Mục tiêu và phạm vi nghiên cứu}
\label{sec:muctieu}
\textbf{Mục tiêu tổng quát:} xây dựng và trình bày một quy trình đánh giá công
bằng, tái lập được cho các thư viện AutoML, đồng thời tiến hành thực nghiệm so
sánh trên dữ liệu thực. Cụ thể, đồ án hướng đến:
\begin{itemize}
    \item Trình bày lại và làm rõ giao thức đánh giá công bằng theo chuẩn
    AMLB~\cite{gijsbers2024amlb}, bao gồm cách kiểm soát fold, ngân sách, độ đo
    và tài nguyên.
    \item Xây dựng hệ thống \textit{AutoML Bench Console} hỗ trợ nạp dữ liệu (CSV,
    OpenML, liên kết Kaggle công khai), tích hợp thư viện AutoML dạng ảnh Docker,
    chạy đánh giá và trực quan hoá kết quả.
    \item Thực nghiệm so sánh FLAML, H2O~AutoML và AutoGluon trên các bộ dữ liệu
    phân lớp và hồi quy, có đối chiếu với các baseline và phân tích các lần thất
    bại.
\end{itemize}

\noindent\textbf{Đối tượng nghiên cứu:} các thư viện AutoML cho dữ liệu dạng bảng
và phương pháp đánh giá so sánh chúng. \\
\textbf{Phạm vi nghiên cứu:} đồ án tập trung vào bài toán học có giám sát trên dữ
liệu dạng bảng (phân lớp và hồi quy); không bao gồm dữ liệu ảnh, văn bản hay
chuỗi thời gian, và không nhằm đề xuất một thuật toán AutoML mới mà tập trung vào
\emph{phương pháp đánh giá} và \emph{công cụ hỗ trợ} tái lập.

\section{Đóng góp của đồ án}
\label{sec:donggop}
Các đóng góp chính của đồ án bao gồm:
\begin{enumerate}
    \item Trình bày lại một cách hệ thống giao thức đánh giá AutoML công bằng theo
    chuẩn AMLB và làm rõ các ``bẫy'' thường gặp khiến so sánh trở nên thiên lệch.
    \item Xây dựng \textit{AutoML Bench Console} --- một hệ thống end-to-end cho
    phép nạp dữ liệu, tích hợp thư viện, chạy đánh giá trong môi trường Docker cô
    lập và trực quan hoá kết quả, với lưu trữ trên PostgreSQL và MinIO.
    \item Thực nghiệm so sánh ba thư viện AutoML tiêu biểu trên nhiều bộ dữ liệu,
    kèm phân tích sự đánh đổi chất lượng -- chi phí và phân tích các lần chạy thất
    bại nhằm bảo đảm tính trung thực khoa học.
\end{enumerate}

\section{Cấu trúc báo cáo}
\label{sec:bocuc}
Phần còn lại của báo cáo được tổ chức như sau.
Chương~\ref{ch:cosoly} trình bày các kiến thức nền tảng về AutoML và các độ đo,
đồng thời khảo sát và phân loại các công trình liên quan để xác định khoảng trống
nghiên cứu.
Chương~\ref{ch:phuongphap} mô tả chi tiết giao thức đánh giá công bằng, kiến trúc
hệ thống, thuật toán quy trình và ví dụ tính toán minh hoạ cách xếp hạng.
Chương~\ref{ch:thucnghiem} trình bày bộ dữ liệu, thiết lập thực nghiệm, kết quả
so sánh, nghiên cứu loại bỏ thành phần (ablation) và phân tích lỗi.
Cuối cùng, Chương~\ref{ch:ketluan} tổng kết đóng góp, nêu giới hạn và đề xuất các
hướng phát triển.
